% UNREGISTERED A.49 candidate; independent audit before inclusion.
\subsection{Complete cover calculations in every proper genus}

\begin{statement}{A.49}{Proved}{A supersingular Kummer family in every genus}
Let \(n\ge2\), choose a prime \(p\equiv3\pmod4\)
with \(p\nmid n\), and let \(a\ge1\). On
\(E:y^2=x^3-x\), choose a point \(Q\) of exact
order \(p\) and a rational function \(f\) with
\[
                         \operatorname{div}(f)=p(O-Q).
                                                               \tag{A.49.1}
\]
All objects are defined over a suitable finite number
field. Let \(C\) be the smooth normalization of
\(z^n=f\), let \(q:C\to E\) be its projection,
and remove the full inverse image of \(E[p^a]\):
\[
                  X=C\setminus q^{-1}(E[p^a]).       \tag{A.49.2}
\]
Then \([p^a]\circ q:X\to U\) is an actual
connected finite etale cover of degree \(np^{2a}\),
with proper genus \(n\) and exactly
\(2+n(p^{2a}-2)\) cusps.

For every coefficient prime, every selected cusp
packet, and every indicated open derived subgroup,
its full geometric pro-prime fixed subgroup is the
selected cusp inertia. For every coefficient prime,
its nontrivial closed procyclic subgroups with open
arithmetic normalizer are exactly the nontrivial
closed subgroups of cusp inertia conjugates.
Thus both assertions of Task A.5 hold in this family
in every proper genus \(n\ge2\), including even
and composite \(n\). The resulting higher-genus
Jacobian need not be assumed to have complex
multiplication.
\end{statement}

\begin{proof}
\paragraph{The actual number-field cover.}
Take a finite extension over which \(Q\), the
\(p^a\)-torsion, and the necessary roots of unity
are defined. Since \([p](O-Q)=0\) in the
elliptic Picard group, the degree-zero divisor
\(p(O-Q)\) is principal. A defining function is
available over a further finite extension; this
justifies the actual rational function in (A.49.1).
The Kummer cover is connected and has degree \(n\),
because the valuation of \(f\) at \(O\) is
\(p\), coprime to \(n\). It is totally ramified
of index \(n\) at \(O,Q\), and unramified
elsewhere. Riemann--Hurwitz gives
\[
                       2g(C)-2=2(n-1),
                      \qquad g(C)=n.                \tag{A.49.3}
\]
In characteristic zero \([p^a]\) is an etale
isogeny of degree \(p^{2a}\). Its origin fibre
is precisely \(E[p^a]\), so removing (A.49.2)
removes all branch points and precisely the origin
fibre of the composite. There is one point over
each of \(O,Q\), and \(n\) points over every
other point of \(E[p^a]\), proving the degree
and cusp count. The Euler characteristic checks
\[
       2-2n-\bigl(2+n(p^{2a}-2)\bigr)=-np^{2a}.
                                                               \tag{A.49.4}
\]
There are infinitely many primes congruent to three
modulo four, by the usual Euclidean proof using
\(4\prod p_i-1\). Hence a suitable \(p\)
exists for every \(n\), and the construction is
not restricted to a congruence class of genera.

\paragraph{The actual model and all its extra cusp sections.}
The displayed elliptic equation has good reduction
at every odd prime. At our chosen \(p\), its
reduction is supersingular: the coefficient criterion
in \cite[Theorem V.4.1(a) and Example V.4.5]{fghSilvermanAEC}
applied to
\(x^{(p-1)/2}(x^2-1)^{(p-1)/2}\) gives zero
coefficient at \(x^{p-1}\), because
\((p-1)/2\) is odd. Thus the special elliptic
curve has \(p\)-rank zero. Every point of
\(E[p^a]\) specializes to its origin: its
specialization is geometric \(p\)-power torsion,
and the special elliptic curve has no nonzero such
points. In particular \(O\) and \(Q\) collide.

Apply the actual two-branch normalization proved in
A.43, with \(d=p\) and \(\gcd(n,p)=1\).
We recall its map, not just its abstract graph.
Starting from the good elliptic model, separate
\(O,Q\) by a chain of blowups. The result consists
of the elliptic main component and a rational chain,
whose last component carries the two distinct marks.
The graph-Laplacian calculation makes the vertical
orders of \(f\) constant. After removing this
common vertical factor and making a finite extension,
its reduction on the main component and the
intermediate chain is a nonzero constant with an
\(n\)-th root. Hensel's lemma, since \(n\) is
invertible in the residue field, splits the actual
normalization there into \(n\) copies.

On the final rational component choose \(u\) with
the branch marks at zero and infinity and the
attaching point at one. Up to a constant the
normalization is
\[
                  z^n=u^p,\qquad u=s^n,
                  \qquad z=s^p.                   \tag{A.49.5}
\]
Bezout's identity for \(n,p\) proves this is the
normalization, and it is a rational component.
There are \(n\) distinct points over \(u=1\).
Thus the actual normalization has \(n\)
supersingular elliptic components, each mapping
isomorphically to the elliptic main component,
joined through rational chains to one rational
centre carrying the two branch marks. The map
\(q\) extends on this normalization to the
blown-up base model, and composing with its
blowdown gives an actual morphism to the original
good elliptic model. Contracting unnecessary chains
is not required in this argument.

It remains to check every extra cusp in (A.49.2),
not only the two branch points. Its image is a
point of \(E[p^a]\), so its special image in
the original good elliptic model is the origin.
On each positive-genus component the extended
map to that model is the isomorphism just described.
The unique point of this component mapping to the
origin is exactly its attaching node with the
rational part of the model. Hence an extra cusp
can specialize only on the rational part or at
one of these attaching nodes. It cannot specialize
at any other smooth point of a positive-genus
component.

After a further finite extension make all cusp
sections rational, and separate them in this
semistable model. These are modifications supported
on the rational part and its attaching nodes.
Separating sections at a smooth rational point
adds rational trees; separating sections at a node
after semistable base change inserts rational chains
and their rational branches. The original elliptic
components persist, and no modification at any other
point of them is needed. Equivalently, those \(n\)
elliptic components already account for all genus
in (A.49.3), so semistable separation cannot introduce
a further positive-genus component or a cycle.
The rational part remains connected, with each
elliptic component attached as a tail. Every cusp
is a smooth mark on that rational part. Thus the
path between any two cusp vertices contains only
rational components, including both endpoints.
This establishes the precise pointed-model hypothesis
for every selected cusp, not merely for the branch
pair.

The proper special Jacobian is the product of these
\(n\) supersingular elliptic Jacobians. It has
good reduction and \(p\)-rank zero. A.41 and
A.42 now give the full fixed-subgroup assertion at
every coefficient prime, for every cusp packet and
every indicated open derived subgroup.

\paragraph{All normalizers, using the same actual component maps.}
Suppose a noncuspidal nontrivial procyclic subgroup
has open arithmetic normalizer. By A.11 its proper
homology has a nonzero generator \(v\), after
restricting the arithmetic field if necessary.
The actual punctured map
\([p^a]\circ q:X\to U\) and A.12 imply that
its proper homology image is zero: a nonzero image
would yield a normalized nontrivial procyclic
subgroup of the target with nonzero proper homology,
contrary to target cuspidality.

If \(\ell\ne p\), the compact-type specialization
gives
\[
                     V_\ell J(C)=\bigoplus_{i=1}^nV_i.
\]
On every component the proper map to the target
elliptic curve is its component isomorphism followed
by \([p^a]\). Thus its map on each \(V_i\)
is a rational Tate-module isomorphism. A nonzero
vector in the total kernel cannot be supported
on a single component. This contradicts the
component-support obstruction of A.47, applied at
residue characteristic \(p\). Its proof is
independent of the residue characteristic and only
requires \(\ell\ne p\).

If \(\ell=p\), the good special Jacobian's
\(p\)-rank zero lets A.46 exclude a noncuspidal
normalized subgroup directly. Conversely every
nonzero closed subgroup of a cusp inertia group
has open arithmetic normalizer after making that
cusp rational. This proves the exact classification.

\paragraph{A concrete even-genus instance.}
Take \(p=3\), \(a=1\), and \(n=4\).
Choose \(Q=(q,r)\) with
\[
          3q^4-6q^2-1=0,
          \qquad r^2=q^3-q,
          \qquad \lambda=\frac{3q^2-1}{2r}.
\]
Then \(Q\) has order three and its tangent
line has divisor \(3Q-3O\). In fact
\[
       x^3-x-\bigl(r+\lambda(x-q)\bigr)^2=(x-q)^3,
                                                               \tag{A.49.6}
\]
using the two displayed relations. One may take
\(f=(y-r-\lambda(x-q))^{-1}\). The smooth
normalization of \(z^4=f\), with the inverse
image of \(E[3]\) removed, has genus four,
\(30\) cusps, and degree \(36\) over \(U\).
Both fixed groups and normalizers have the asserted
classification at every coefficient prime. The
accompanying exact certificate verifies this tangent
identity, the ramification and Euler counts, and
the tame normalization exponents for sample parameters;
the preceding proof applies to every \(n,p,a\)
in the statement.
\end{proof}
